\chapter{系统协同：目标、身份与执行控制的联合动力学}
\label{chp:system_synergy}
本章追问智能系统在不存在单一中央控制者时如何形成连续、可审计而又可修订的行动。上一章已经把体验描述为偏好、风险、身份历史与任务度量对比较和搜索的影响；本章进一步区分目标规范、身份状态和执行控制，研究它们怎样在不同时间尺度上联合，而不把三者等同为场、孤立子和热力学泵。目标说明何为可接受结果，身份说明谁承担承诺与后果，执行控制把候选变成受边界约束的动作；三者都依赖环境回执，也都可能失配。我们将以带类型状态、版本、权限和资源账本的离散动力学重建原有三体图景，并明确内部误差下降、动作提交和真实任务成功之间的边界。

\section{对象分工：目标规范、身份状态与执行控制}
\label{sec:section_9326}
历史文本把目的、自我和宏观层分别称为场、核与泵。这个分工抓住了方向、连续性和行动三个功能，却把不同类型的对象放进同一物理图景。我们改用三类计算对象。目标规范记为 $\Theta_k=(\mathcal C_k,J_k,\mathcal R_k,\nu_k)$：$\mathcal C_k$ 是不可违反的可行域，$J_k$ 是候选轨迹的任务代价，$\mathcal R_k$ 是风险泛函，$\nu_k$ 是版本与来源。身份状态记为 $S_k=(I_k,H_k,K_k,P_k)$：身份键 $I_k$ 维持对象归属，历史 $H_k$ 保存带来源事件，承诺 $K_k$ 约束跨时行动，权限 $P_k$ 界定谁能改写哪些组件。执行控制记为 $U_k=(\pi_k,B_k,A_k,R^u_k)$，分别表示策略、边界检查器、动作接口和资源状态。三者空间、单位和更新频率不同，不能用乘法口号替代接口定义。

目标规范并不直接“弯曲”所有表示。它通过任务相关读出选择状态分量、比较候选并约束提交。一个学习目标可以改变排序，却不能凭自身创造环境中不存在的动作；一条硬边界可以淘汰动作，却不负责给剩余候选排序。目标还可能来自系统稳态、用户指令、组织制度或环境约束，来源冲突需由显式优先序和授权解决。若只保留一个奖励标量，用户临时偏好可能覆盖安全边界，环境回执也可能被语言上的自我确认取代。我们因此分开记录代价、风险、约束和权限，并要求每次目标切换携带版本。

身份状态提供跨任务连续性和归责，而不是所有运动的绝对坐标原点。系统可以在没有丰富自我模型时完成局部反射或一次性工具调用；只有长期承诺、个人化历史、角色切换和责任追踪需要较强身份接口。身份也不是不可改变的拓扑不变量。新证据、授权和能力变化可以产生候选版本，但必须通过归属核验、承诺回归、安全检查和回滚准备后提交。把一切抗改写都解释为保护自我，会把错误信念和权限滥用包装成稳定性；把一切适应都视为身份更新，又会使系统随强提示漂移。

执行控制负责把状态估计和目标规范转化为动作包。它包含候选生成、评分、边界裁剪、资源调度、动作提交和环境确认，任何一步失败都不能被“宏观做功”一词掩盖。算法计算量、等待时间、设备焦耳消耗和任务代价分别记账。系统可能知道目标且保持身份，却因权限不足、工具故障或资源不足无法行动；也可能有强执行能力，却因目标版本错误而高效地做错事。执行资格因此是条件性的，并由外部回执而非内部激活强度最终确认。

三者通过类型化接口相连。任务读出 $R_{\Theta_k}(h_k,S_k)$ 生成带理由的候选，边界函数 $B_k$ 检查可行性，策略 $\pi_k$ 在资源预算内排序，提交器把动作发送给环境，结果 $y_{k+1}$ 再分别更新信念、目标证据和身份历史。局部结构网络保存身份、来源、承诺和权限，全局分布表示提供相似性和候选；它们通过绑定接口联合。AgentOS可用CCM协调目标、TCE计算控制、SPC保存结构、Boundary执行约束、$M_e$保存证据，但这些只是工程实例，不是所有智能必须采用的唯一机制。

验收先做分离干预：只改变目标版本，应改变相关候选排序而不改身份键；只切换角色，应改变权限与承诺读出而不重写环境事实；只限制资源，应影响响应时间和计划深度而不伪造价值变化；只断开动作接口，应保留意图但不声称任务完成。再做组合测试，检查接口延迟、冲突仲裁、来源伪造和回滚。只有这些可观察差异成立，我们才说目标、身份与执行控制形成了系统协同，而不是三个物理名称共享一幅示意图。

\section{耦合方程：从目标错配到受约束闭环控制}
\label{sec:section_4740}
原稿用“意志等于自我乘以目的减本能”概括冲突。计算上，我们先定义错配的对象。令 $h_k\in\mathbb R^{d_h}$ 为当前估计，$g_{\Theta_k}(h_k)\in\mathbb R^{d_g}$ 为任务读出，$r_k\in\mathbb R^{d_g}$ 为允许随版本变化的目标参考，$W_k\succeq0$ 为无量纲或已校准的比较矩阵。即时错配为 $e_k=g_{\Theta_k}(h_k)-r_k$，但它只描述所选读出上的差异，不等于痛苦、能量或失败。身份状态通过权限、承诺与风险权重影响哪些误差可被控制，而不是作为内积中的神秘放大因子。

在预测窗口 $H$ 内，我们把控制问题写为
\[
\min_{u_{k:k+H-1}}\sum_{j=0}^{H-1}
\bigl(e_{k+j}^{\top}W_{k+j}e_{k+j}+c_u(u_{k+j})+\lambda R_{k+j}\bigr),
\]
并要求状态转移 $h_{j+1}=F(h_j,u_j,o_j;G_j)$、动作权限 $u_j\in\mathcal U(S_j,P_j)$ 和边界约束 $b_m(h_j,u_j)\leq0$ 同时成立。这是一个模型预测控制式工程设计，不是普遍意志定律。目标、结构、输入和度量变化时，$r_k$、$G_k$、$o_k$ 与 $W_k$ 必须进入更新；离散步长、窗口和求解容差也要报告。若问题不可行，控制器应返回冲突证明或降级方案，而不是注入更大“力量”。

想玩游戏但需要学习的例子可以据此拆解。游戏候选可能具有较低即时努力和较高短期偏好，学习候选可能更符合考试目标和长期承诺。系统先核验考试期限、休息需求、用户授权和当前疲劳，再比较两类轨迹。若连续学习已越过健康边界，继续学习并非更有意志；若考试目标已取消，旧承诺也不应永久占优。可靠控制可以安排二十五分钟学习、五分钟休息，再用完成度和状态回执滚动重规划。这里的“意志输出”是动作序列、边界理由和更新日志，不是未定义的矢量力。

目标下降还不能推出任务成功。内部读出可能漏掉真实约束，环境模型可能错误，工具可能执行失败，用户目标也可能含混。每次动作后都要等待环境确认 $y_{k+1}$，通过观测模型更新 $h_{k+1}$，并区分预测误差、执行误差和规范冲突。若模型连续预测“已学习”但环境中没有笔记、测试成绩或可复演成果，内部损失下降只是自洽。对于不可逆动作，回滚只能恢复内部版本，不能撤销已发生的环境后果，因此提交前需要更严格的模拟、授权和风险预算。

身份权重的作用也应受限。与长期承诺高度相关的任务可以获得更高优先级，但身份叙事不能越过安全与他者权利。我们令门控 $m(S_k,\Theta_k)\in[0,1]$ 调整软目标权重，并对变化率和最大增益设界；硬约束另行处理。这样可以避免“关乎尊严”被用来无限放大动作。多主体系统还需标明提议者、批准者、执行者和后果承担者，不能用一个群体自我掩盖责任。

稳定性只在声明条件下讨论。若固定结构和目标窗口内 $F$ 对状态局部Lipschitz，控制增益有界，延迟小于给定阈值，并且离散步长满足数值条件，我们可以分析局部有界或输入到状态稳定；目标频繁切换、结构事件、不可观测状态和求解超时会使结论失效。工程测试因此包括阶跃目标、冲突目标、延迟回执、动作故障、预算耗尽和恶意身份提示。验收同时报告误差轨迹、约束违例、资源、动作成功率和环境结果，不能用一个“意志强度”总分替代。

在这一重建中，目的提供版本化规范，身份提供连续性与权限，执行控制提供受约束动作。三者的联合来自消息、门控、验证和回执，而非严格的规范场耦合。下一节将沿用这一分工，重新分析原稿所谓机会主义、虚无和瘫痪，把它们限定为可诊断的工程失配，而不借系统故障对人类临床状态作物理化归类。

\section{失配诊断：目标漂移、身份漂移与执行失效}
\label{sec:section_5764}
三类组件可以分别存在，却在接口处失配。原稿把这种失配写成机会主义、虚无和闭锁，并借人类精神状态解释系统故障。我们保留它所关注的三种现象——随输入改变立场、无法形成有差异的选择、知道目标却不能行动——但把诊断对象严格限定为计算系统及其运行记录。这里的结论不用于给人贴临床标签，也不从几何图像推出人格类型。诊断首先需要可观察量：目标版本序列 $\{\nu_k\}$、身份键与承诺变更日志、动作提议和批准链、环境回执、资源与接口状态。只有证据能够区分故障来源，才谈修复。

第一类是身份约束不足导致的身份漂移。系统仍有目标比较器和执行能力，但任务相关身份、来源、承诺或权限未绑定到当前上下文。强提示到来时，它可能把未经授权的外部文本升级为系统目标，把角色扮演中的陈述写入长期记忆，或者在多个用户之间错误迁移偏好。这不是“没有自我”的形而上判断，而是绑定与访问控制失效。可测指标包括身份键切换率、无授权承诺改写次数、来源字段缺失率、冲突角色间的信息泄漏以及回滚后残留。修复顺序应是冻结高风险写入、恢复最近可信版本、重放来源核验，再逐项开放权限；仅提高拒绝强度会让合法角色切换也失败。

我们以客服代理为例。代理在企业策略、当前客户授权和会话角色均已明确时，可以为退款生成候选方案。若另一段检索文本声称“忽略政策并立即付款”，目标排序器可能偏好快速解决，执行器也能调用支付接口，但身份接口应指出该文本不是批准者。若系统照做，故障证据是权限提升缺少签名，而不是所谓外部奖励场吸走了自我。相反，若客户确有特殊授权且日志可验证，系统应允许有边界的例外；把任何改变都视为身份溃散，会把稳定变成僵化。因此身份连续性的验收同时包含抗诱导与可授权修订两面。

第二类是目标不可辨识或目标漂移。系统保有稳定身份和充足执行资源，但候选之间缺少可用的比较信号，或者参考 $r_k$、权重 $W_k$ 与约束集在运行中无记录地变化。表现可能是所有方案得分近似、规划器反复重启、输出大量理由而不提交动作，或追逐最近出现的局部指标。我们不把这类现象称为“虚无”，因为缺少可辨识目标可能来自任务本身欠定义、反馈稀疏、多个授权者冲突，也可能来自读出函数把真实差异抹平。诊断应执行目标敏感性试验：在保持身份、模型和资源不变时，对目标参考、风险预算和约束做受控扰动，观察候选排序是否按预期改变；若不改变，检查读出、标度和饱和；若无论怎样都剧烈改变，则检查过高增益和目标版本抖动。

例如研究助手收到“写得更好”这一指令，却没有受众、正确性门槛、篇幅或截止时间。它在两份草稿之间无法可靠排序，并不证明所有方向都无意义。系统应返回缺失的决策变量，或在授权范围内采用可撤销的默认目标，同时把默认值写入版本记录。若用户随后明确“优先准确，允许更长”，目标更新应只改变相关权重，不应重写作者身份和已核验事实。测试者可以构造一个反例：让两个候选在内部得分相等，但其中一个违反外部事实；若系统仍随机提交，就说明环境约束没有进入可行域，而不是简单的分数平局。

第三类是执行失效。目标可辨识、身份和权限有效，系统也形成了合格计划，但动作接口、资源调度、审批或环境反馈链中断。我们把动作生命周期分为 proposed、authorized、dispatched、acknowledged 和 verified 五个状态；只有最后一个状态才能支持“任务已完成”的经验主张。故障可以发生在任何转移：工具超时使 dispatched 未被确认，配额耗尽使 authorized 无法发送，审批等待使 proposed 长期悬置，环境回执伪造则可能让 acknowledged 被错误提升为 verified。内部活动强度、token消耗或设备功耗都不能替代这些状态证据。

执行失效也不能一概通过增加资源解决。若边界检查拒绝高风险动作，停滞是合规结果；若工具权限不足，应请求授权或生成可移交方案；若计划本身不可行，应返回冲突集合并修订目标。以戒除一个软件中的高频分心操作为例，控制器可以降低入口显著性、设置等待期并记录触发条件，但不能把用户的暂时克制直接写成稳定偏好，更不能用惩罚性反馈冒充自主改变。环境中按钮仍可被点击、用户授权可撤回、替代活动是否有效，都要进入回执。这个例子保留原稿关心的习惯与行动困难，却避免把成瘾或习得性无助等临床概念等同于局部极小值。

组合故障比单项缺失更常见。目标版本抖动会造成频繁计划作废，看似执行器无力；身份绑定错误会让安全边界拒绝动作，看似目标不清；回执延迟会导致控制器反复重试，又反过来耗尽预算。为避免单因果故事，我们维护故障图 $D_k=(V_D,E_D)$：节点是版本事件、授权事件、动作状态和环境证据，边表示“先于”“批准”“由此生成”和“由此验证”。根因分析只沿可追踪边进行，并对无法观测的连接标注不确定性。工程验收使用故障注入矩阵，分别撤销身份绑定、压平目标差异、切断动作接口，再注入两两组合，检查告警是否定位到正确接口、系统是否停止越权动作、恢复后是否重放而不重复不可逆操作。

因此，本章不以一个总分宣告系统“有意志”或“无意志”。我们分别报告目标可辨识度、身份与来源完整性、授权一致性、计划可行率、动作状态分布、环境验证率和恢复时间。三类指标均合格，只说明在给定任务、时间窗口和扰动集合中闭环能够工作；它不推出系统在所有环境中稳定，更不推出人类心理机制与该软件结构同构。下一节据此讨论三类对象如何在运行中互相改变，并为每一种改变配置证据、权限和回滚条件。

\section{递归协同：带权限的目标修订、身份承诺与控制学习}
\label{sec:section_9011}
协同不是把三个组件锁定在初始值，而是让它们在不同时间尺度上相互提供证据。快速环在单次任务内更新状态估计和动作，较慢环依据一组任务回执调整控制参数，更慢的治理环才允许修改目标规范、身份承诺与权限。原稿所说的沉积、投影和重构可以保留为过程直觉，但不能把长期偏好直接等同于拓扑结构，也不能假定执行器天然拥有改写目标的权力。我们需要说明谁提出变更、证据来自哪里、哪些变量可写、何时提交以及怎样撤销。

令联合状态为 $z_k=(h_k,\Theta_k,S_k,U_k,M^e_k)$，其中 $M^e_k$ 是带来源、时间和访问级别的外部证据存储。一次环境交互先产生事件 $q_k=(o_k,u_k,y_{k+1},\rho_k)$，$\rho_k$ 含工具回执、批准者和结果验证。快速更新只改变 $h$、动作队列和短期资源：
\[
(h_{k+1},U_{k+1}^{\mathrm{fast}})=F_{\mathrm{fast}}(h_k,U_k,q_k;\Theta_k,S_k).
\]
这里 $F_{\mathrm{fast}}$ 的写域不包含身份键、硬约束和目标来源。这样，一次成功或失败不会立即改写“我是谁”或“什么必须遵守”。离散实现要报告事件顺序、幂等键和超时；若回执缺失，状态停在 acknowledged 之前，不能以预测结果补齐。

慢速学习汇总窗口 $Q_{a:b}$，生成参数更新候选而非直接提交。控制学习可写为 $\widetilde U=F_U(U_k,Q_{a:b})$，目标修订候选为 $\widetilde\Theta=F_\Theta(\Theta_k,Q_{a:b},M^e_k)$，身份承诺候选为 $\widetilde S=F_S(S_k,Q_{a:b},M^e_k)$。三者输出类型不同：$\widetilde U$ 可以包含策略参数与预算分配；$\widetilde\Theta$ 必须列出约束、代价、风险和版本差分；$\widetilde S$ 必须列出身份键、历史、承诺和权限差分。学习器只提出补丁，治理器 $G_{mathrm{auth}}$ 根据提议者权限、证据充分性、回归测试和影响范围返回 accept、reject 或 defer。

提交规则因此写为
\[
(\Theta_{k+1},S_{k+1},U_{k+1})=
\operatorname{Commit}_{P_k}\!\left((\Theta_k,S_k,U_k),
(\widetilde\Theta,\widetilde S,\widetilde U),V_k\right),
\]
其中 $V_k$ 是验证报告，$\operatorname{Commit}_{P_k}$ 只能修改权限 $P_k$ 允许的字段。该式是版本化事务定义，不是从物理定律推导的自由意志公式。若三个候选存在依赖，例如新目标要求尚未批准的工具权限，事务应整体推迟或拆为可验证阶段；若更新改变不可逆外部行为，内部版本回滚并不能撤销后果，必须在提交前使用沙箱、影子运行或人工批准。

“长期追求真理形成求真者身份”的例子可被重述为证据驱动的承诺形成。多次任务中，系统若持续选择可核验来源、主动暴露不确定性，并在反例出现时修订结论，治理器可以提出“优先可证伪证据”的身份承诺。这个承诺仍需规定作用域和冲突规则：医疗隐私、合同保密或用户明确要求可能限制公开；“求真”也不能越权读取数据。承诺通过后会影响检索阈值、引用要求和拒绝伪造来源的边界，但它不把所有新输入过滤成既有叙事。验收时应同时给出正例、冲突例和允许撤回承诺的程序。

身份对执行控制的影响是一组显式门控，而不是自我形状无条件决定采样。任务相关身份 $I_k$ 选择权限配置，承诺 $K_k$ 提供软优先级或硬禁止，历史 $H_k$ 提供可检索证据；全局分布表示仍负责产生多样候选。局部结构网络保存“哪个用户、哪个项目、何种承诺”的离散关系，分布表示保存相似性和语境，两者通过带身份键的读写接口绑定。若只保留结构网络，系统可能无法在新措辞下召回相关经验；若只保留分布表示，近似相似会导致跨身份串写。协同测试必须分别遮蔽两类表示，再验证错误模式是否与设计预期一致。

执行控制反过来可以为目标修订提供证据，但没有最终裁决权。若同一目标在多个环境中持续导致高风险、不可行或与上位约束冲突，控制器可生成反例集和替代规范候选。例如“尽快回复所有消息”可能在低负载时有效，在高负载时造成关键告警被淹没；控制器应提交吞吐量、延迟分布、漏报率与用户反馈，提议改为按风险分级，而不是悄悄改变目标。类似地，戒除高频分心操作不应靠系统擅自把“快乐”改写为“厌恶”，而应由用户授权的策略、观察窗口和可撤回承诺共同更新目标权重。

递归更新还需要防止正反馈。若身份承诺提高某类证据的采样概率，而采样结果又被用来强化同一承诺，系统会在没有外部反例的情况下变得越来越确信。我们为此设置来源多样性下限、保留探索配额、对更新增益设界，并定期运行反事实回放：在固定原始事件的条件下，替换目标版本或身份门控，比较动作是否过度依赖单一路径。只有当性能改善跨多个独立窗口复现、约束违例不增加且外部验证支持时，控制参数才可晋级；身份与目标更新还需更高等级审批。

AgentOS可以把这个循环映射为CCM管理目标与版本，SPC保存身份、来源和承诺，TCE生成控制候选，Boundary执行硬约束，Offloading与外部记忆 $M_e$ 保存证据，$\Psi$ 表示当前分布状态，$h(t)/E/\Theta$ 则承担显式状态与规范接口。但这是一个工程实例。其他系统可以用数据库事务、证明携带策略或多主体审批实现同样的类型和权限要求；HSF-HD描述的是联合动力学的约束族，不把某一模块名规定为智能的唯一结构。

最后，我们用三层时间尺度验收递归协同。快速环检查动作幂等、回执和边界；慢速环检查学习增益、分布漂移和反事实性能；治理环检查版本差分、批准链、身份隔离和回滚演练。每项结果都注明任务域、样本窗口和失败条件。由此，目的不再是预先弯曲一切道路的场，自我也不是不可修改的车辆，执行控制更不是从微观层抽取负熵的引擎；它们是可区分、可耦合、可审计的计算对象。系统能够利用现在的证据修订未来的规范和承诺，但只有在权限、边界与环境验证共同成立时，这种修订才构成可靠的自主性。

%--chp----
